%-------------------------------------------------------%
%-                          包                         -%
%-------------------------------------------------------%
\usepackage{CJK,CJKnumb}
\usepackage{url}
\usepackage[backend=biber,style=numeric-comp,sorting=none,gbpub=false,bibstyle=gb7714-2015,citestyle=gb7714-2015,
gbpunctin=false,gbnamefmt=lowercase,
]{biblatex}

\usepackage{amsmath, amsfonts, amssymb}
\usepackage{booktabs} %三线表
\usepackage{threeparttable}
\usepackage{multicol}
\usepackage{multirow}
\usepackage{fancyhdr}
\usepackage{latexsym}
\usepackage{mathrsfs}
\usepackage{wasysym}
\usepackage{enumerate}
\usepackage{titlesec}
\usepackage{epstopdf}
\usepackage{caption}
\usepackage{graphicx, subfig} %图包
\usepackage{float}
\usepackage{setspace}
\usepackage{array}
\usepackage{fancyhdr} %页眉页脚包
\usepackage{fontspec} %英文字体包
\usepackage{titletoc} %设置目录格式
\usepackage{geometry} %设置页边距
% \usepackage{circledsteps}%enumerate带圈的序号
\usepackage{pifont}
\usepackage{adjustbox}


\usepackage[labelsep=quad]{caption}
\captionsetup[figure]{font=footnotesize, font=bf}
\captionsetup[table]{font=footnotesize, font=bf}

% \usepackage{enumitem}
% \usepackage{float}%提供float浮动环境
% \usepackage{booktabs}%提供命令\toprule、\midrule、\bottomrule
% \usepackage{graphicx}




%论文尺寸为A4，全文页面设置为：上下2.54cm，左右3.17cm，含有大量图表可微调。页边距0.5cm
\geometry{a4paper,left=3.17cm,right=3.17cm,top=2.54cm,bottom=2.54cm,bindingoffset=0.25cm} 

\setCJKmainfont[Path=/System/Library/Fonts/Supplemental/]{Songti.ttc}
% \setCJKmainfont{simsun.ttc} %文中中文为宋体
\setmainfont{Times New Roman} %文中英文为新罗马
\setlength{\baselineskip}{20pt} %行距为固定值20磅


%---------------------------------------------------------------------------%
%-                        字体设置                               -%
%---------------------------------------------------------------------------%
\setCJKfamilyfont{song}[Path=/System/Library/Fonts/Supplemental/]{Songti.ttc}
\setCJKfamilyfont{fs}[Path=/System/Library/Fonts/Supplemental/]{Songti.ttc}
\setCJKfamilyfont{kai}[Path=/System/Library/Fonts/Supplemental/]{Songti.ttc}
\setCJKfamilyfont{hei}[Path=/System/Library/Fonts/]{STHeiti Medium.ttc}
\setCJKmonofont[Path=/System/Library/Fonts/Supplemental/]{Songti.ttc}

\newcommand{\song}{\CJKfamily{song}}    % 宋体   (simsun.ttc)
\newcommand{\fs}{\CJKfamily{fs}}        %仿宋体  (simfs.ttf)
\newcommand{\kai}{\CJKfamily{kai}}      % 楷体   (simkai.ttf)
\newcommand{\hei}{\CJKfamily{hei}}      % 黑体   (simhei.ttf)

%设置加粗体
%加粗楷体
\setCJKfamilyfont{kaitib}[Path=/System/Library/Fonts/Supplemental/,AutoFakeBold]{Songti.ttc}
\newcommand{\kaitib}{\CJKfamily{kaitib}}
%加粗仿宋体
\setCJKfamilyfont{fangsongti}[Path=/System/Library/Fonts/Supplemental/,AutoFakeBold]{Songti.ttc}
\newcommand{\fangsongti}{\CJKfamily{fangsongti}}
%加粗黑体
\setCJKfamilyfont{heitib}[Path=/System/Library/Fonts/,AutoFakeBold]{STHeiti Medium.ttc}
\newcommand{\heitib}{\CJKfamily{heitib}}
%加粗宋体
\setCJKfamilyfont{songtib}[Path=/System/Library/Fonts/Supplemental/,AutoFakeBold]{Songti.ttc}
\newcommand{\songtib}{\CJKfamily{songtib}}
%正文加粗快捷键
\newcommand{\tbf}{\bfseries \songtib}

%---------------------------------------------------------------------------%
%-                          字号设置                            -%
%---------------------------------------------------------------------------%

% \newcommand{\chuhao}{\fontsize{42pt}{\baselineskip}\selectfont}
% \newcommand{\xiaochuhao}{\fontsize{36pt}{\baselineskip}\selectfont}
% \newcommand{\yihao}{\fontsize{26pt}{\baselineskip}\selectfont}
% \newcommand{\xiaoyihao}{\fontsize{24pt}{\baselineskip}\selectfont}
% \newcommand{\erhao}{\fontsize{22pt}{\baselineskip}\selectfont}
% \newcommand{\xiaoerhao}{\fontsize{18pt}{\baselineskip}\selectfont}
\newcommand{\sanhao}{\fontsize{16pt}{\baselineskip}\selectfont}
\newcommand{\xiaosanhao}{\fontsize{15pt}{\baselineskip}\selectfont}
\newcommand{\sihao}{\fontsize{14pt}{\baselineskip}\selectfont}
\newcommand{\daxiaosihao}{\fontsize{12pt}{\baselineskip}\selectfont} %%%%%%% by Hanlin Cai
\newcommand{\xiaosihao}{\fontsize{11pt}{\baselineskip}\selectfont} %%%%%%% by Hanlin Cai
\newcommand{\wuhao}{\fontsize{10.5pt}{\baselineskip}\selectfont}
\newcommand{\xiaowuhao}{\fontsize{9pt}{\baselineskip}\selectfont}
\newcommand{\liuhao}{\fontsize{7.5pt}{\baselineskip}\selectfont}
\newcommand{\qihao}{\fontsize{5.5pt}{\baselineskip}\selectfont}

%---------------------------------------------------------------------------%
%-                               Chinesization                             -%
%---------------------------------------------------------------------------%
\newtheorem{theorem}{\hskip 2em定理}[section]
\newtheorem{definition}{\hskip 2em定义}[section]
\newtheorem{exam}{\hskip 2em例}[section]
\newtheorem{note}{\hskip 2em注}[section]
\newtheorem{proof}{{\it \hskip 2em\textbf{证明}}}[section]
\newtheorem{solution}{{\it \hskip 2em\textbf{解}}}
\renewcommand{\tablename}{Table}
\renewcommand{\figurename}{Figure}
\renewcommand{\refname}{\centerline {Reference}}
\renewcommand{\contentsname}{\heitib\sanhao 目\quad 录}
\renewcommand{\thefootnote}{\arabic{footnote}}
\renewcommand{\theequation}{\thesection-\arabic{equation}}
\renewcommand{\thefigure}{\thesection-\arabic{figure}}
\renewcommand{\thetable}{\thesection-\arabic{table}}

%---------------------------------------------------------------------------%
%-                     公式随章编号                               -%
%---------------------------------------------------------------------------%
\date{}
\makeatletter
\renewcommand*\l@chapter[2]{%
 \ifnum \c@tocdepth >\m@ne
   \addpenalty{-\@highpenalty}%
   \vskip 1.0em \@plus\p@
   \setlength\@tempdima{1.5em}%
   \begingroup
     \parindent \z@ \rightskip \@pnumwidth
     \parfillskip -\@pnumwidth
     \leavevmode \bfseries
     \advance\leftskip\@tempdima
     \hskip -\leftskip
     #1\nobreak\leaders\hbox{$\m@th
       \mkern \@dotsep mu\hbox{-}\mkern \@dotsep
       mu$}\hfill \nobreak\hb@xt@\@pnumwidth{\hss #2}\par
     \penalty\@highpenalty
   \endgroup
 \fi}

% \date{}
% \numberwithin{equation}{section}
% \numberwithin{figure}{section}
% \numberwithin{table}{section}

%---------------------------------------------------------------------------%
%-                              Tabular Option                             -%
%---------------------------------------------------------------------------%
\newcommand{\tabincell}[2]{\begin{tabular}{@{}#1@{}}#2\end{tabular}}

%---------------------------------------------------------------------------%
%-                              Caption Option                             -%
%---------------------------------------------------------------------------%
\captionsetup[figure]{labelsep=space}
\captionsetup[table]{labelsep=space}

%------------------------------------------------------%
%-                    文献来源                          -%
%-------------------------------------------------------%
\addbibresource{Text/reference.bib}

%------------------------------------------------------%
%-                Page Control 1                      -%
%------------------------------------------------------%
\topmargin=0cm \oddsidemargin=0.5cm \textwidth=15cm
\textheight=22cm
\setlength{\headheight}{13pt}

\titlecontents{section}[0em]{\song\wuhao\addvspace{3pt}}{\contentslabel{2.5em}}{}%
{\titlerule*[0.7pc]{$\cdot$}\contentspage}%

\titlecontents{subsection}[2em]{\song\wuhao}{\contentslabel{3em}}{}%
{\titlerule*[0.7pc]{$\cdot$}\contentspage}%

\titlecontents{subsubsection}[4em]{\song\wuhao}{\contentslabel{3.5em}}{}%
{\titlerule*[0.7pc]{$\cdot$}\contentspage}%


\renewcommand{\baselinestretch}{1.1}


%------------------------------------------------------%
%-                Page Control 2                      -% Hanlin Cai
%------------------------------------------------------%

\titleformat{\section}[display]{\centering\fontsize{16pt}{12pt}\selectfont\bfseries}{}{0pt}{}

\titleformat{\subsection}
{\fontsize{14pt}{12pt}\selectfont\bfseries}{\thesubsection}{1em}{}

\titleformat{\subsubsection}
{\fontsize{13pt}{12pt}\selectfont\bfseries}{\thesubsubsection}{1em}{}

\titleformat{\subsubsection}
{\fontsize{12pt}{12pt}\selectfont\bfseries}{\thesubsubsection}{1em}{}
